% Generated draft: requires mathematical and source review.
\documentclass[11pt]{article}
\usepackage[margin=25mm]{geometry}
\usepackage{fontspec}
\setmainfont{Latin Modern Roman}
\usepackage{amsmath,amssymb,mathtools,unicode-math}
\setmathfont{Latin Modern Math}
\usepackage{enumitem,needspace,xurl,hyperref}
\setlength{\parindent}{0pt}
\setlength{\parskip}{5pt}
\setlength{\emergencystretch}{3em}
\newcommand{\N}{\mathbb{N}}
\newcommand{\Z}{\mathbb{Z}}
\newcommand{\Q}{\mathbb{Q}}
\newcommand{\R}{\mathbb{R}}
\newcommand{\C}{\mathbb{C}}
\newcommand{\F}{\mathbb{F}}
\newcommand{\A}{\mathbb{A}}
\newcommand{\PP}{\mathbb{P}}
\newcommand{\E}{\mathbb{E}}
\newcommand{\eps}{\varepsilon}
\newcommand{\dd}{\,\mathrm{d}}
\begin{document}
\section*{UnsolvedMath ID: 800975}
\textbf{Source problem code:} catalogue entry
\subsection{Definitions and mathematical statement}
A closed surface is a compact surface without boundary. Using the usual orientable-genus convention, let \(S\) be a connected, closed, orientable smooth surface of genus \(g\geq 2\), and let \(h\) be a smooth Riemannian metric on \(S\). Its geodesic flow is represented on the cotangent bundle \(T^*S\), equipped with its canonical symplectic form, by the Hamiltonian
\[
H_h(q,p)=\frac12 h_q^{-1}(p,p).
\]
A first integral is a function constant along this Hamiltonian flow; for a smooth function \(F\), this is equivalent to \(\{H_h,F\}=0\), where \(\{\cdot,\cdot\}\) is the canonical Poisson bracket.

The supplied question, attributed to Bangert, asks: does there exist a closed surface of genus at least \(2\) admitting a Riemannian metric whose geodesic flow is completely integrable?

One conventional smooth Liouville interpretation gives the following explicit existence question:
\[
\begin{gathered}
\exists\, S,h,F,\qquad \operatorname{genus}(S)\geq 2,\qquad F\in C^\infty(T^*S,\R),\\
\{H_h,F\}=0,\qquad
\dd H_h\wedge\dd F\neq 0
\quad\text{on a dense open subset of }T^*S\setminus 0_S,
\end{gathered}
\]
where \(0_S\) denotes the zero section and \(S,h\) satisfy the conditions above. In this interpretation, the two commuting integrals \(H_h\) and \(F\) are functionally independent on that dense open set.

\emph{Unresolved qualification.} The supplied excerpt does not define complete integrability. In particular, it does not specify the regularity and domain of the additional integrals, the required extent of functional independence, or conditions on singular invariant sets. The displayed formulation is therefore a conditional interpretation, not an established identification of the source's intended definition. The excerpt also leaves the connectedness and orientability conventions implicit.
\subsection{Short English statement}
Bangert's question asks whether some closed surface of genus at least \(2\) admits a Riemannian metric with completely integrable geodesic flow. A standard smooth Liouville formulation uses an additional first integral independent of the geodesic Hamiltonian on a dense open set, but the supplied excerpt does not specify whether this is the intended meaning. The supplied registry labels the problem open; this status has not been independently verified.
\subsection{Sources}
\begin{itemize}
\item Local UnsolvedMath registry, permanent ID \(800975\), source problem code \texttt{catalogue entry}, in the AMR Open Problem Lists collection. The supplied record has upstream ID \(800975\) and upstream revision \texttt{d371ee917ac9647fbb06cd6fc529deec6f3dda82}.
\item Burns and Matveev, \emph{Open Problems and Questions about Geodesics}, listed in the supplied record as 2019, Question 56. Supplied source URL: \url{https://arxiv.org/abs/1308.5417}. The record attributes the question/context to Bangert and lists Burns and Matveev in its proposer field.
\item The supplied provenance records access on 2026-10-16 and extraction by \texttt{arxiv\_latex}. These are supplied provenance details, not a claim of independent source access or verification. The registry's supplied status is \texttt{open}.
\end{itemize}
\end{document}
